\begin{abstract}
无法把遥测发送给托管模型的组织，只能用小型开源权重 LLM 实现告警分诊自动化，而这些模型恰恰在调查流程上失败：探测却不收敛、不给出结论，或者服从写进日志的指令。我们研究应当从模型手里拿走哪些分诊决定，以及由此形成怎样的安全边界。\hesp{} 是一个控制器：它维护候选原因的账本，按单位成本的期望信息增益选择只读探针，只接受有当前证据支持的结论，并且可以自行结束调查。在一项基于三个合成任务族的受控研究中（七项采用内部冻结协议的研究，五个 7B 到 72B 的模型，21{,}278 个经过审计的 LLM 回合），“探测什么”与“何时停止”是两种独立的失败。采用验证器的证据标准后，控制器无需 LLM 就能完成每一个案例；预先登记的比较没有发现让 Qwen2.5-72B 决定何时停止带来增益，而即使给基线同样清楚的提示，它们仍然远远落后。要求一条决定性观测是以自动化换取安全：联合证据能解决更多“原因使用相同工具”的案例，但会把一部分攻击结为与之相似的合法活动。结束守卫在每个模型上都挡住了“关闭告警”的注入指令；伪造的来源能击败控制器，除非良性结论需要来自不同来源的佐证。模型剩下的作用是读取原始日志：小模型能读懂让规则解析器失效的新格式，但可能采纳攻击者声称的读数。
\end{abstract}
